\chapter{Estado Actual del Conocimiento (Antecedentes)}
\label{ch:antecedentes}

Este capítulo sistematiza las contribuciones de la literatura relevante en ocho ejes. Los primeros cuatro provienen de la literatura argentina reciente: la insuficiencia de los enfoques contables \citep{rodriguez2023, sanches2019}, el análisis de la economía política del endeudamiento y las reestructuraciones \citep{aruguete2021, roldan2021}, las consecuencias materiales del ciclo financiero sobre la estructura productiva \citep{bohoslavsky2024, cantamutto2024}, y la literatura que cuestiona la urgencia de la fatiga fiscal en contextos de tasa de interés real baja \citep{blanchard1990, blanchard2019}. Los cuatro ejes restantes provienen de la literatura internacional de simulación estocástica de deuda y de dos antecedentes específicos sobre Argentina que, sin integrar ambas tradiciones, dan origen directo al diseño metodológico del Capítulo \ref{ch:metodologia}: el linaje VAR más simulación estocástica del que desciende el SVAR restringido de esta tesis \citep{medeiros2012, adler2014, wijayanti2020}, las dos lecturas convergentes de la antesala de la crisis de 2001--2002 \citep{cetrangolo1997, calvo2003}, la controversia entre el ratio Deuda/PIB y el balance patrimonial completo \citep{levyyeyati2021}, y el costo real \textit{ex post} del evento de \textit{default} \citep{hebert2017}.

\section{Primer Eje: Limitaciones de los Enfoques Contables Lineales}
\label{sec:eje1}

La literatura coincide en señalar las limitaciones de las metodologías deterministas clásicas para evaluar la solvencia en economías periféricas.

\subsection{La perspectiva integral de \texorpdfstring{\textcite{rodriguez2023}}{Rodríguez (2023)}}
Desde el marco de la Nueva Economía Keynesiana, \textcite{rodriguez2023} analiza la evolución de la deuda pública argentina entre 2001 y 2022 mediante un indicador de ``efectos sobre la deuda pública'' que descompone la dinámica del cociente Deuda/PIB en sus determinantes. Su principal hallazgo es que, durante el 68\% del período relevado, dicho indicador se ubicó en terreno negativo, evidenciando señales de insostenibilidad e incapacidad de maniobra fiscal, y que el cociente Deuda/PIB tomado de manera aislada resulta insuficiente para extraer conclusiones sobre sostenibilidad.

\subsection{El canal cambiario del perfil de deuda: \texorpdfstring{\textcite{sanches2019}}{Sanches (2019)}}
En su tesis de maestría, \textcite{sanches2019} descompone el cociente Deuda/PIB argentino en sus determinantes para aislar el peso específico del tipo de cambio, dado el perfil de deuda con elevada participación en moneda extranjera, y complementa el análisis con un estudio de panel entre países emergentes sobre la sensibilidad del diferencial de rendimiento soberano al perfil de moneda de la deuda. Su hallazgo central es que, para un perfil de deuda como el argentino, una depreciación real afecta el cociente Deuda/PIB tanto por el canal del capital como por el de los intereses.

\section{Segundo Eje: Procesos de Reestructuración y Actores Internacionales}

\subsection{Crisis de balanza de pagos y reestructuración: \texorpdfstring{\textcite{aruguete2021}}{Aruguete (2021)}}
\citet{aruguete2021} examina la crisis de balanza de pagos que Argentina transitó en 2020, originada en el proceso de sobreendeudamiento iniciado cuatro años antes, y el proceso de reestructuración de la deuda en moneda extranjera, habilitado formalmente por la declaración de emergencia pública en materia económica, financiera, fiscal y de deuda de la Ley 27.541 \citep{ley27541}, que buscó restaurar la sostenibilidad mejorando el perfil de vencimientos y despejando las necesidades de divisas de corto plazo. Su diagnóstico subraya que los modelos de sostenibilidad no pueden obviar el uso macroeconómico de los fondos prestados ni la vulnerabilidad que genera un perfil de vencimientos concentrado en moneda extranjera.

\subsection{Teoría de juegos y Cláusulas de Acción Colectiva: \texorpdfstring{\textcite{roldan2021}}{Roldán (2021)}}
En su tesina de grado, \citet{roldan2021} analiza la renegociación de 2020 de la deuda soberana argentina con los tenedores privados de instrumentos bajo jurisdicción extranjera, abordando el proceso desde la Teoría de Juegos y destacando el rol de las Cláusulas de Acción Colectiva, incorporadas a los nuevos instrumentos tras la experiencia litigiosa de la crisis de 2014, para obligar a los acreedores a alinearse una vez superado cierto umbral de aceptación, reduciendo la capacidad de presión de los acreedores litigantes (\textit{holdouts}).

\section{Tercer Eje: La Tensión Material y la Restricción Externa}

\subsection{Sostenibilidad de la deuda en sentido amplio: \texorpdfstring{\textcite{bohoslavsky2024}}{Bohoslavsky y Cantamutto (2024)} y \texorpdfstring{\textcite{cantamutto2024}}{Cantamutto et al. (2024)}}
\citet{bohoslavsky2024} sostienen que el sobreendeudamiento argentino guarda una relación estructural con la especialización productiva extractiva del país, de modo que un criterio de sostenibilidad restringido al balance fiscal corriente omite la dimensión estructural de largo plazo. En la misma línea, \citet{cantamutto2024} proponen leer la sostenibilidad desde la perspectiva de la ``sostenibilidad de la vida'', cuestionando que el criterio fiscal, centrado en la capacidad de pago frente a los acreedores, pueda evaluarse con independencia de sus efectos sobre la reproducción social de la población; el planteo retoma y sistematiza un argumento formulado antes en \citet{feliz2021}, sobre la tensión entre la sostenibilidad de la deuda y la sostenibilidad de la vida en el caso argentino.

\section{Cuarto Eje: Literatura Crítica sobre la Urgencia de la Fatiga Fiscal}

Junto a los tres ejes anteriores, que dan por plausible la fatiga fiscal, corresponde incorporar la literatura que cuestiona su urgencia. Ya \citet{blanchard1990} propuso matizar el cociente Deuda/PIB con indicadores fiscales complementarios en lugar de tratarlo como límite autosuficiente; tres décadas después, \textcite{blanchard2019} argumenta que en contextos donde la tasa de interés real se ubica sistemáticamente por debajo del crecimiento económico ($r<g$), el gobierno puede sostener déficits primarios moderados sin que el cociente diverja, lo cual relativiza la urgencia atribuida a los superávits primarios crecientes. Desde una crítica metodológica distinta, \citet{everaert2017} advierten que, en estimaciones panel de la función de reacción fiscal, lo que se interpreta como fatiga fiscal puede confundirse con heterogeneidad de parámetros entre países agrupados bajo un mismo coeficiente de reacción; esa advertencia refuerza la conveniencia de un diseño de series de tiempo de caso único, como el de esta tesis, frente a un panel que asuma homogeneidad de la reacción fiscal entre economías estructuralmente distintas. Es la misma advertencia que recibe, desde el eje de la simulación estocástica de deuda, el hallazgo original de \citet{medeiros2012} que da origen al marco institucional hoy aplicado a Argentina (Sección \ref{sec:eje_simulacion}): un coeficiente de fatiga fiscal estimado sobre un panel heterogéneo de países no es directamente transferible a un caso único sin volver a contrastarlo.

Esta perspectiva refuerza la decisión de someter la hipótesis de fatiga fiscal a un contraste explícito de significatividad estadística, en lugar de asumirla por analogía con economías centrales con acceso irrestricto a financiamiento en moneda propia, supuesto que el Marco Teórico (Capítulo \ref{ch:marco_teorico}) descarta para el caso argentino.

\section{Quinto Eje: Simulación Estocástica de Deuda y Marcos de Riesgo Soberano}
\label{sec:eje_simulacion}

Un quinto eje, de origen internacional y no específico de Argentina, aporta el linaje metodológico directo del protocolo econométrico de esta tesis (Capítulo \ref{ch:metodologia}).

\subsection{De la simulación estocástica de deuda a los marcos de riesgo soberano}
La literatura empírica sobre sostenibilidad de deuda soberana se desplazó, en la última década y media, de las proyecciones deterministas de un único escenario base hacia ejercicios de simulación estocástica capaces de generar distribuciones de probabilidad sobre la trayectoria de la deuda. El punto de partida empírico de este giro es \citet{medeiros2012}, quien estimó un Vector Autorregresivo (VAR) irrestricto y una Función de Reacción Fiscal de panel para quince países de la Unión Europea, combinando ambos instrumentos para simular miles de trayectorias de Deuda/PIB y construir los primeros \textit{fan charts}\footnote{Representación gráfica de una distribución de probabilidad de trayectorias futuras de una variable, mediante bandas de percentiles que se abren como un abanico a medida que aumenta el horizonte de proyección.} de uso sistemático en un organismo oficial (la Comisión Europea). Ese hallazgo fue incorporado después como supuesto de comportamiento en marcos de sostenibilidad de organismos multilaterales, entre ellos el Marco de Riesgo Soberano y Sostenibilidad de la Deuda del FMI (SRDSF), aplicado a Argentina en su Octava Revisión del Servicio Ampliado \citep{imf2024}.

\subsection{De la macroeconomía general a las economías emergentes exportadoras de commodities}
Un segundo cuerpo de literatura empírica trasladó el problema de sostenibilidad a economías emergentes expuestas a shocks externos de términos de intercambio y a interrupciones súbitas del financiamiento. \citet{adler2014}, dentro del volumen editado por el FMI sobre América Latina, estimaron modelos VAR globales (GVAR) y de panel (PVAR) para un conjunto de países de la región, introduciendo como variable exógena internacional el Índice de Precios Netos de Materias Primas (NCPI) desarrollado por \citet{adlermagud2014}, que pondera cada \textit{commodity} por la posición neta exportadora o importadora del país. Su resultado central, que más de la mitad de la variabilidad de las entradas de capital de corto plazo a economías como la argentina responde a condiciones internacionales (tasa de la Reserva Federal, índice VIX) y no a variables de atracción local, refuerza la idea de que la sostenibilidad de la deuda en la región depende en gran medida de factores \textit{push}\footnote{Determinante externo al país receptor que impulsa movimientos de capital, en contraste con los factores \textit{pull}, asociados a condiciones domésticas de atracción de inversión.} ajenos a la política doméstica; el mismo estudio documenta, sin embargo, una asimetría específica de Argentina, donde el capital residente tiende a sumarse a la salida de capitales en lugar de repatriarse para amortiguar una crisis, como ocurre en economías desarrolladas comparables.

En paralelo, la literatura de sistemas de alerta temprana (\textit{early warning systems})\footnote{Modelos estadísticos diseñados para anticipar la probabilidad de una crisis mediante indicadores de alerta temprana, sin pretender explicar el mecanismo causal subyacente.} buscó anticipar crisis de deuda soberana con una estrategia metodológica distinta: no explicar mecanismos, sino maximizar poder predictivo. \citet{wijayanti2020} aplicaron el ratio de ruido a señal de Kaminsky, Lizondo y Reinhart y una regresión logística binomial sobre un panel de 43 países en desarrollo entre 1960 y 2017, prediciendo el 61,5\% de las crisis de deuda con dos años de anticipación bajo un umbral de probabilidad del 30\%. Ese resultado, moderado pero no trivial, expone una tensión relevante para cualquier antecedente apoyado en el marco del FMI: los sistemas de alerta temprana maximizan poder predictivo agregado sobre un panel amplio de países, mientras que el \textit{fan chart} estocástico del SRDSF busca representar la distribución de riesgo de un país específico; ambos abordajes pueden, y de hecho suelen, dar señales distintas frente a un mismo episodio.

\section{Sexto Eje: Dos Lecturas Convergentes de la Crisis de 2001--2002}

Dentro de la literatura específicamente centrada en Argentina, dos trabajos independientes, con instrumentos metodológicos casi opuestos, llegaron a diagnósticos convergentes en la antesala de la crisis de 2001--2002. \citet{cetrangolo1997} reconstruyeron el esquema ahorro-inversión del Sector Público Nacional no Financiero para 1975--1996, desagregaron el presupuesto en ocho rubros de ingresos y siete de gasto, vincularon cada uno a su variable macroeconómica relevante mediante coeficientes fijos, y aplicaron el filtro de Hodrick-Prescott para separar tendencia de ciclo y calcular un déficit macroeconómicamente ajustado y un impulso fiscal siguiendo la metodología de \citet{blanchard1990}. Su conclusión fiscal fue tranquilizadora: el riesgo fiscal cíclico bajo Convertibilidad era pequeño (0,3--0,6\% del PIB), muy inferior a la volatilidad de los años setenta y ochenta. Sin embargo, el propio documento advirtió, en su escenario central de proyección para 1997--2001, que el desequilibrio de cuenta corriente externa crecía hasta superar el 7\% del PIB hacia 2001, anticipando, desde una lente estrictamente fiscal, que el verdadero riesgo no era fiscal sino de sostenibilidad externa.

\citet{calvo2003}, que trabajaron con un instrumental completamente distinto (contabilidad de balanza de pagos, calibración de elasticidades de demanda de no transables y la identidad estándar de dinámica de deuda), llegaron a la misma advertencia por otra vía. Su marco de economía cerrada, dolarizada y con descalce de balances mostró que, ante una interrupción súbita del financiamiento (\textit{sudden stop})\footnote{Interrupción abrupta e inesperada del financiamiento externo hacia una economía, término acuñado en la literatura de crisis financieras internacionales.} como la desatada por la crisis rusa de 1998, la magnitud del ajuste de tipo de cambio real necesario para restablecer el equilibrio externo dependía críticamente de la apertura comercial y del descalce de moneda de los pasivos, dos rasgos en los que Argentina se ubicaba en el peor extremo posible entre los casos comparados (Argentina, Brasil, Chile, Colombia, Ecuador). El diálogo entre ambos trabajos es directo: uno detecta el problema por el lado del flujo (cuenta corriente), el otro por el lado del stock (descalce patrimonial), y ambos, escritos antes o inmediatamente después del \textit{default}, coinciden en que el diagnóstico fiscal de corto plazo no bastaba para juzgar la sostenibilidad.

\section{Séptimo Eje: El Balance Patrimonial Completo frente al Ratio Deuda/PIB}

Dos décadas después, la literatura sobre Argentina volvió sobre un problema metodológico contiguo al del Primer Eje: si el ratio Deuda/PIB es o no una medida confiable de riesgo cuando media un salto cambiario abrupto. El propio \citet{rodriguez2023} (Sección \ref{sec:eje1}) atribuye la caída del ratio Deuda/PIB durante 2003--2011 a un efecto contable de apreciación del tipo de cambio real antes que a un desendeudamiento estructural; el mismo mecanismo, en sentido inverso, explica por qué el salto cambiario de fines de 2023 llevó el ratio al 156,7\% del PIB y disparó, dentro del SRDSF del FMI, una alarma de ``Riesgo Alto'' (ancho de \textit{fan chart} de 122,7 puntos) que el propio staff del organismo debió matizar cualitativamente por no reflejar el riesgo real de mediano plazo \citep{imf2024}.

\citet{levyyeyati2021} llevaron esta discusión un paso más allá al proponer reemplazar el ratio Deuda/PIB por el patrimonio neto del sector público (valor presente de impuestos menos valor presente de gastos y pasivos explícitos), simulado mediante \textit{bootstrap} de un VAR de tres variables (PIB, tipo de cambio real, tasa internacional). Aplicado a Argentina, el ejercicio arrojó un patrimonio neto medio de 2,2 veces el PIB sin territorio negativo en la distribución simulada, y un hallazgo llamativo: una devaluación real del 20\% apenas mueve ese patrimonio neto (de 2,225 a 2,214 del PIB), porque el mayor costo en dólares de la deuda se compensa con más recaudación de rubros transables y la licuación de gastos en pesos. Este resultado entra en tensión directa con el cálculo tradicional que los propios autores citan de \citet{calvo2003}: bajo el enfoque clásico de descalce de moneda, una devaluación del 50\% eleva el ratio Deuda/PIB argentino de 36,5\% a 50,8\%. Dos antecedentes centrados en el mismo país y el mismo tipo de shock llegan a conclusiones opuestas sobre su severidad fiscal según si el instrumento utilizado es el ratio de deuda o el balance patrimonial completo; los propios autores advierten que su resultado benigno depende de que Argentina, tras 2001, redujo fuertemente la dolarización de su deuda pública, y que el mismo ejercicio aplicado a la estructura de deuda de 2001 habría arrojado un deterioro fiscal severo ante el mismo shock.

\section{Octavo Eje: El Costo Real del Evento de \textit{Default}}

Un último antecedente no discute la sostenibilidad \textit{ex ante}, sino el costo \textit{ex post} del evento que estos marcos buscan evitar. \citet{hebert2017} explotaron el litigio República Argentina contra NML Capital como experimento natural: los fallos judiciales sucesivos alteraron exógenamente la probabilidad de \textit{default} percibida por el mercado, sin depender de las noticias económicas del día. Mediante identificación por heterocedasticidad, estimaron que un aumento del 10\% en la probabilidad de \textit{default} (medida por CDS) provocó una caída del 6\% en el índice ponderado de acciones argentinas, y que el paso de 40\% a 100\% de probabilidad de \textit{default}, lo efectivamente observado en 2014, implicó una destrucción de valor de mercado de las firmas cotizantes de entre 30\% y 45\%, concentrada en exportadoras y filiales de multinacionales. Este antecedente complementa, desde el lado de las consecuencias reales, a la literatura de sostenibilidad \textit{ex ante} revisada arriba: cuantifica el costo del evento de cola que los marcos estocásticos (\citealp{medeiros2012}; \citealp{imf2024}) intentan mantener con baja probabilidad, y que el enfoque de balance completo de \citet{levyyeyati2021} sugiere que, bajo la estructura de deuda actual, sería menos probable que en 2001.

\section[Área de Vacancia: la Endogeneización Pendiente del Riesgo País]{Área de Vacancia: la Endogeneización Pendiente del Riesgo País}
\label{sec:vacancia}

Delimitado el estado del arte en sus ocho ejes, corresponde precisar la vacancia científica que esta tesis busca cubrir. La literatura local relevada, tanto \citet{rodriguez2023} y \citet{sanches2019} desde el eje contable-estocástico, como \citet{aruguete2021} y \citet{cantamutto2024} desde la economía política, incluso en su vertiente con referato \citep{bohoslavsky2024}, coincide en tratar el esfuerzo fiscal y el riesgo soberano como fenómenos paralelos y narrativamente vinculados, pero no sometidos a un contraste de causalidad simultánea explícito: se argumenta cualitativamente que el riesgo país erosiona el espacio fiscal, pero no se estima un modelo que instrumente el EMBI+ para aislar su efecto causal sobre el resultado primario, purgado de la causalidad inversa (un peor resultado fiscal eleva el riesgo percibido por los acreedores). La literatura de origen internacional (Quinto Eje) resuelve esa endogeneidad solo a nivel de panel agregado, no para un caso único; y ni ella ni el enfoque de balance completo de \citet{levyyeyati2021} (Séptimo Eje) llegan a instrumentar el riesgo soberano dentro de un sistema que además lo trate como variable conjuntamente endógena con el resultado fiscal y el tipo de cambio. Esta vacancia no es exclusiva de la literatura local: \citet{gonzalez2023}, en la evaluación más reciente de funciones de reacción fiscal para el conjunto de América Latina y el Caribe, ni siquiera logran incluir a Argentina entre los diez países con series suficientemente limpias (sin \textit{default} ni hiperinflación extrema) para una estimación individual por país, y el coeficiente positivo que reportan para el grupo en el que Argentina queda agrupada surge de una estimación de panel agrupado, no de una estimación específica para el caso argentino. Esta tesis atiende precisamente esa laguna, mediante un diseño de serie de tiempo de caso único que incorpora explícitamente los quiebres estructurales de la economía argentina en lugar de promediarlos con los de otros países de la región.

Esta omisión resulta central: si el riesgo país se trata como exógeno en una regresión simple, cualquier coeficiente estimado se encuentra potencialmente sesgado por simultaneidad. La vacancia científica de esta tesis radica en la \textbf{endogeneización sistémica del Riesgo País (EMBI+)}, tratado como variable conjuntamente endógena dentro del VECM y del SVAR restringido del Capítulo \ref{ch:metodologia} y corregido puntualmente mediante Variables Instrumentales en Dos Etapas (IV-2SLS) sobre la ventana original, y en el contraste explícito, mediante remuestreo estocástico, de significatividad estadística de cualquier umbral de fatiga fiscal detectado, lo que permite evaluar estadísticamente los planteos de la literatura previa (Capítulo \ref{ch:resultados}).

La Tabla \ref{tab:contraste_estado_arte} sintetiza esta vacancia en forma de matriz de contraste, cruzando las tres dimensiones metodológicas centrales del protocolo econométrico de esta tesis (tratamiento del riesgo país, modelización de no linealidades y naturaleza de las proyecciones) contra el tratamiento que reciben en cada uno de los dos ejes de la literatura relevada.

\begin{table}[H]
\centering
\caption{\label{tab:contraste_estado_arte} \textit{Matriz de contraste del Estado del Arte y posicionamiento de la Tesis.}}
\begin{tabular}{p{3.1cm}p{3.6cm}p{3.6cm}p{3.9cm}}
\toprule
\textbf{Dimensión / Autor} & \textbf{Enfoque Contable (Rodríguez, Sanches)} & \textbf{Economía Política (Aruguete, Cantamutto)} & \textbf{Contribución de esta Tesis} \\
\midrule
Tratamiento del Riesgo País & Tratado narrativamente como variable de contexto o variable de aproximación de costo de financiamiento & Tratado narrativamente como canal de vulnerabilidad externa (Aruguete) y subordinado a sus costos sociales (Cantamutto) & Endogeneizado explícitamente mediante instrumentación IV-2SLS (VIX, ETF EMB) \\
Modelización de No Linealidades & Asumida cualitativamente (``bola de nieve'', escenarios discretos) sin contraste formal & Asumida cualitativamente (rechazo político y barrera material como límites conjuntos) sin contraste formal & Testeada formalmente mediante modelo de umbral de Hansen con significatividad por remuestreo (Sup-LM) \\
Naturaleza de las Proyecciones & Deterministas contables o perturbaciones aisladas sobre escenarios puntuales & Ausentes o subordinadas a la narrativa histórica e institucional, sin proyección cuantitativa & Análisis de Sostenibilidad de Deuda (DSA) estocástico con simulación de Monte Carlo (1.000 iteraciones) y gráficos de abanico de probabilidad \\
\bottomrule
\end{tabular}
\begin{flushleft}
\footnotesize \textit{Nota.} Elaboración propia en base a la revisión de \citet{rodriguez2023}, \citet{sanches2019}, \citet{aruguete2021}, \citet{bohoslavsky2024} y \citet{cantamutto2024} sistematizada en las Secciones anteriores de este capítulo.
\end{flushleft}
\end{table}

La Tabla \ref{tab:contraste_estado_arte} delimita la contribución de esta tesis: el diferencial no reside en la extensión muestral, pues ambos ejes disponen de series históricamente tan extensas, sino en el refinamiento de la identificación causal en presencia de sesgos por simultaneidad. Mientras la literatura previa infiere el rol del riesgo país y de la fatiga fiscal a partir de la coincidencia narrativa entre series, esta tesis somete ambos mecanismos a un protocolo de identificación explícito (instrumentación del EMBI+, contraste de significatividad del umbral de Hansen) que distingue entre lo que la evidencia disponible permite afirmar y lo que continúa siendo una hipótesis interpretativa plausible. El aporte principal es de carácter metodológico: reemplaza la asociación narrativa de riesgo soberano y esfuerzo fiscal por un diseño de identificación causal que aísla el componente exógeno de cada mecanismo propuesto por la literatura precedente.
